
\section{Complete study-11 metrics}\label{app:basefull}
Table~\ref{tab:basefull} gives every metric for every model of the Veltri
2018 study, including AUPRC, which the main text omits for space. The two
ensemble rows used the tune partition for their thresholds (0.56 average,
0.575 stack); base models use the default 0.5.

\begin{table}[h]\centering\small
\caption{All study-11 models on the Veltri test partition (percentages
except MCC).}
\label{tab:basefull}
\begin{tabular}{lcccccc}
\hline
Model & SENS & SPEC & ACC & MCC & AUC & AUPRC \\
\hline
PepCNN seed 7 & 89.47 & 90.73 & 90.10 & 0.8020 & 96.04 & 96.33 \\
PepCNN seed 27 & 94.66 & 80.48 & 87.57 & 0.7591 & 95.67 & 95.95 \\
PepCNN seed 77 & 89.04 & 91.43 & 90.24 & 0.8050 & 96.01 & 96.32 \\
PepGNN seed 11 & 76.26 & 86.24 & 81.25 & 0.6281 & 88.14 & 87.48 \\
PepGNN seed 111 & 87.08 & 73.31 & 80.20 & 0.6097 & 87.93 & 86.37 \\
RF (descriptors) & 90.31 & 88.76 & 89.54 & 0.7908 & 96.19 & 96.62 \\
Average ensemble & 90.87 & 90.45 & 90.66 & 0.8132 & 95.95 & 96.38 \\
Stacked ensemble & 90.31 & 92.13 & 91.22 & 0.8246 & 96.61 & 97.03 \\
\hline
\end{tabular}
\end{table}

The AUPRC column reinforces the ranking: the stack (97.03) leads every
base model, and the GNN seeds trail on every metric, consistent with the
capacity and pooling argument of Appendix~\ref{app:arch}. Seed variance
among the CNNs is small on AUC (95.67--96.04) but larger on MCC
(0.759--0.805), the expected pattern when only the threshold varies.
